\section{文献综述与研究假设}

\subsection{企业数字化转型与 AI 应用研究}

数字化转型通常被界定为企业借助数字技术改造业务流程、组织结构和价值创造方式的系统性过程 \citep{Vial_2019,Verhoef_2021}。相关研究强调，数字化转型并非简单的信息化投入，而是涉及组织能力更新、管理方式重构和战略再定位的持续过程 \citep{Warner_2019}。在这一过程中，人工智能、大数据分析、平台化系统和自动化工具被视为企业重塑经营效率和创新能力的重要抓手。

不过，数字化转型研究也指出，企业“谈论数字化”与“真正完成数字化能力建设”之间存在明显差距。企业可能出于战略沟通、市场预期管理或对外展示需要，强化对 AI 和数字化的表述，但这些表述并不必然对应已经落地的流程再造、资源投入和组织支撑。换言之，企业在披露文本中使用数字化叙事，既可能是能力信号，也可能只是概念占位。因此，仅凭是否提及 AI 或数字化，并不足以判断企业数字化转型的真实性。

\subsection{文本披露、信息质量与叙事包装研究}

信息披露研究普遍认为，文本叙事是企业与外部投资者之间缓解信息不对称的重要机制 \citep{Healy_2001,Verrecchia_1990}。随着文本分析方法的发展，越来越多研究从可读性、语气、语义结构和叙事风格角度衡量企业披露质量 \citep{LOUGHRAN_2016,LOUGHRAN_2014,Li_2008}。相关结论表明，文本特征不只是“表达形式”，往往还承载着关于企业治理质量、信息透明度与未来表现的额外信息。

进一步地，印象管理研究指出，企业管理层可能通过夸大措辞、模糊边界、选择性叙事和事后合理化等方式影响外部受众认知 \citep{MerklDaviesConcept2011,MerklDaviesSensemaking2011,Henry_2008}。语气过度积极、抽象概念堆砌和承诺化表述，可能提高文本的说服力，却未必提升其信息含量。尤其当企业处于强激励、强展示的情境时，叙事包装更可能出现。IPO 正是这样一种典型情境：企业需要同时面向投资者和监管者讲清楚自己的成长逻辑，因此文本既可能传递真实信息，也可能存在包装性修辞。

基于上述研究，可以把企业的 AI 披露区分为两个面向。其一是真实性，即文本是否体现为有信息量、可验证、具体且边界清晰的能力信号；其二是包装性，即文本是否更多依赖宏大叙事、口号化修辞和未来承诺。二者并非简单对立，但在经验上往往呈现出不同的信息含量和兑现可能性。

\subsection{IPO 文本披露与问询回复研究}

IPO 场景的一个重要特征，是招股说明书集中展示企业商业模式、技术能力、风险因素与募投安排，因而成为高信息密度文本。已有研究发现，IPO 招股书的语言特征和信息含量能够帮助识别企业质量，并影响发行定价和后续表现 \citep{Hanley_2010,Loughran_2013}。与上市后年报相比，IPO 文本更集中地体现企业如何在资本市场首次讲述自己的成长故事。

在中国注册制背景下，招股书之外的问询与回复进一步提供了一个有价值的治理窗口。问询制度强化了监管者对披露充分性、业务模式、核心技术和成长逻辑的追问，使企业需要对原始叙事进行补充说明、澄清和辩护。这意味着，若企业在招股书中对 AI 与数字化进行高强度表述，但在问询中遭遇对真实性、先进性或必要性的持续追问，那么其数字化叙事的可信度就值得进一步考察。相反，若企业不仅在招股书中给出清晰场景，而且在问询回复中能够提供投入、系统、流程和团队等客观证据，则更可能说明其数字化表述具有较高真实性。

现有研究对 IPO 文本已有较多讨论，但针对“AI 叙事真实性”这一更具体问题的研究仍然较少，特别是将招股说明书与问询回复联合起来识别“真实信号还是叙事包装”的研究更为有限。本文试图在这一缺口上推进。

\subsection{理论分析与研究假设}

从信号理论视角看，企业在 IPO 阶段披露 AI 与数字化信息，本质上是在向外部市场传递能力信号 \citep{Spence_1973}。如果这类表述能够落到具体业务场景，伴随研发投入、系统建设、技术团队和实施进度等证据，并在问询回复中保持解释一致，那么其更可能反映真实能力基础。此类企业具备更高概率将数字化投入转化为经营效率、成长能力和组织协同优势，因此上市后经营表现也更可能改善。据此，提出假设：

\textbf{H1：}企业 IPO 阶段 AI 相关披露真实性越高，上市后经营表现越好。

从印象管理视角看，若企业在披露中大量使用抽象战略表述、宏大口号和承诺化措辞，却缺少具体业务场景、客观投入证据和边界说明，那么这类文本更接近叙事包装而非真实能力信号 \citep{MerklDaviesConcept2011,Huang2014}。包装性叙事虽然可能在短期内增强文本吸引力，但由于其信息含量不足，往往更难在上市后转化为持续经营改善。若问询过程中监管者还持续追问其合理性、先进性或披露充分性，则进一步表明这类叙事的可兑现性较弱。据此，提出假设：

\textbf{H2：}企业 IPO 阶段 AI 相关披露的叙事包装性越强，上市后经营表现越弱。
